% ── §1.5: Συνοπτικά Αποτελέσματα ─────────────────────────────

\section{Συνοπτικά Αποτελέσματα}
\label{sec:results_preview}

Η υλοποίηση του συστήματος FL είναι ολοκληρωμένη και λειτουργική. Τα
αρχικά πειραματικά αποτελέσματα στο σύνολο δεδομένων AAL αναδεικνύουν ότι:

\begin{itemize}
    \item Το προτεινόμενο 1D-CNN, εκπαιδευμένο ομοσπονδιακά (FedAvg) σε
    22~ετερογενείς (Non-IID) χρήστες για 20~γύρους, επιτυγχάνει
    $\mathbf{76{,}8\%}$ ακρίβεια (F1-macro $= 0{,}768$) σε \emph{ανεξάρτητο}
    (held-out) σύνολο ελέγχου.

    \item Σε σύγκριση με κεντρικοποιημένη εκπαίδευση ($84{,}8\%$, άνω
    φράγμα), η απώλεια των $\mathbf{8{,}0}$ ποσοστιαίων μονάδων
    ποσοτικοποιεί το πραγματικό κόστος της ομοσπονδιακής μάθησης υπό
    Non-IID ετερογένεια --- το αντίτιμο για τη διατήρηση της
    ιδιωτικότητας~\cite{Li2020}.

    \item Η ανάλυση αποτίμησης συνεισφοράς (σε εξέλιξη) δείχνει ήδη ότι ο
    \emph{όγκος} δεδομένων ενός χρήστη δεν προβλέπει την πραγματική του
    συνεισφορά: οι χρήστες με τα περισσότερα δεδομένα δεν συμπίπτουν με
    εκείνους που βελτιώνουν περισσότερο το καθολικό μοντέλο (κατά
    Leave-One-Out). Το εύρημα αυτό θεμελιώνει την ανάγκη για μηχανισμό
    κινήτρων βασισμένο στην \emph{ποιότητα} και όχι στην ποσότητα της
    συνεισφοράς.
\end{itemize}

\noindent
Η οριστική ποσοτική σύγκριση των μετρικών συνεισφοράς (Leave-One-Out
έναντι Monte-Carlo Shapley) και η τελική επιλογή του μηχανισμού κινήτρων
βρίσκονται υπό ολοκλήρωση και αποτελούν το επόμενο στάδιο της εργασίας.
